\documentclass[12pt]{article}
\usepackage{amsmath,amssymb,amsthm}
\usepackage[margin=1in]{geometry}

\newtheorem{theorem}{Theorem}
\newtheorem{lemma}[theorem]{Lemma}

\title{Problem 195: Inscribed Circles of Triangles\\with One Angle of 60 Degrees}
\author{Project Euler}
\date{}

\begin{document}
\maketitle

\section{Problem Statement}

How many triangles with integer sides having one angle of exactly $60^\circ$ have an incircle with radius not exceeding $R = 10^5$?

\section{Mathematical Foundation}

\begin{theorem}[Incircle Radius]
For a triangle with sides $a$, $b$ and included angle $60^\circ$, the opposite side is $c = \sqrt{a^2 + b^2 - ab}$, the area is $\Delta = \frac{\sqrt{3}}{4}\,ab$, and the incircle radius is
\[
r = \frac{\sqrt{3}\,ab}{2(a + b + c)}.
\]
\end{theorem}

\begin{proof}
By the law of cosines, $c^2 = a^2 + b^2 - 2ab\cos 60^\circ = a^2 + b^2 - ab$. Area: $\Delta = \frac{1}{2}ab\sin 60^\circ = \frac{\sqrt{3}}{4}ab$. Incircle: $r = \Delta/s = \frac{\sqrt{3}\,ab}{2(a+b+c)}$. \qed
\end{proof}

\begin{theorem}[Eisenstein Integer Norm]
The form $a^2 - ab + b^2$ is the norm $N(a + b\omega)$ in the Eisenstein integers $\mathbb{Z}[\omega]$, $\omega = e^{2\pi i/3}$. This ring is a UFD (Euclidean domain).
\end{theorem}

\begin{proof}
Standard: the Euclidean function is the norm $N(z) = |z|^2 = x^2 - xy + y^2 \geq 0$, and $N(z) = 0$ iff $z = 0$. \qed
\end{proof}

\begin{theorem}[Primitive Triple Parameterization]
All primitive triples $(a, b, c)$ with $c^2 = a^2 - ab + b^2$, $\gcd(a,b) = 1$, fall into two families parameterized by coprime $m > n > 0$:

\textbf{Family 1} ($m \not\equiv n \pmod{3}$):
\[
a = 2mn + n^2, \quad b = m^2 - n^2, \quad c = m^2 - mn + n^2.
\]

\textbf{Family 2} ($m \equiv n \pmod{3}$):
\[
a = \frac{2mn + n^2}{3}, \quad b = \frac{m^2 - n^2}{3}, \quad c = \frac{m^2 - mn + n^2}{3}.
\]
\end{theorem}

\begin{proof}
Factor $a^2 - ab + b^2 = (a - b\omega)(a - b\bar{\omega})$ in $\mathbb{Z}[\omega]$. For this to be a perfect square, write $a - b\omega = \epsilon \alpha^2$ for unit $\epsilon$ and $\alpha = m + n\omega$. Expanding $\alpha^2$ and matching components yields both families. When $m \equiv n \pmod{3}$, all components are divisible by 3, giving Family~2 after reduction. \qed
\end{proof}

\begin{lemma}[Scaling]
For primitive triple $(a_0, b_0, c_0)$ with incircle radius $r_0$, the triple $(ka_0, kb_0, kc_0)$ has radius $kr_0$. Valid copies: $\lfloor R/r_0 \rfloor$.
\end{lemma}

\begin{proof}
$r(ka_0, kb_0, kc_0) = \frac{\sqrt{3}\,k^2 a_0 b_0}{2k(a_0+b_0+c_0)} = kr_0$. \qed
\end{proof}

\begin{lemma}[Multiplicity]
Each unordered pair $\{a, b\}$ with $a \neq b$ gives one triangle. The parameterization counts $(a,b)$ and $(b,a)$ separately, so each such primitive triple contributes $2\lfloor R/r_0 \rfloor$ to the total.
\end{lemma}

\begin{proof}
Swapping $a$ and $b$ preserves $c$ and $r$. The 60-degree angle is uniquely determined. \qed
\end{proof}

\section{Algorithm}

\begin{verbatim}
for each coprime (m, n) with m > n > 0:
    Compute (a, b, c) for Family 1 or 2
    Compute r0 = sqrt(3)*a*b / (2*(a+b+c))
    Add floor(R/r0) * multiplicity to count
\end{verbatim}

\section{Complexity Analysis}

\begin{itemize}
    \item \textbf{Time:} $O(B^2)$ where $B = O(R)$ is the upper bound on $m$. Most pairs are pruned by coprimality and the $r_0 \leq R$ bound.
    \item \textbf{Space:} $O(1)$.
\end{itemize}

\section{Answer}

\[
\boxed{75085391}
\]

\end{document}
